\documentclass[11pt,a4paper]{article}

% ---------- Packages ----------
\usepackage[utf8]{inputenc}
\usepackage[T1]{fontenc}
\usepackage{lmodern}
\usepackage[margin=0.9in]{geometry}
\usepackage{microtype}
\usepackage{parskip}
\usepackage{enumitem}
\usepackage{xcolor}
\usepackage{titlesec}
\usepackage{hyperref}
\usepackage{fancyvrb}
\usepackage{listings}
\usepackage{tcolorbox}
\usepackage{fontawesome5}
\usepackage{textcomp}
\usepackage{tabularx}

% ---------- Colors ----------
\definecolor{primary}{HTML}{2563EB}
\definecolor{secondary}{HTML}{4B5563}
\definecolor{codebg}{HTML}{F3F4F6}
\definecolor{codefg}{HTML}{1F2937}
\definecolor{rulec}{HTML}{D1D5DB}

% ---------- Hyperref ----------
\hypersetup{
    colorlinks=true,
    linkcolor=primary,
    urlcolor=primary,
    citecolor=primary,
    pdftitle={FinTrack Technical Report},
    pdfauthor={Sagar Makwana}
}

% ---------- Section formatting ----------
\newcommand{\sectionrule}{\vspace{2pt}\titlerule[0.6pt]}
\titleformat{\section}
  {\Large\bfseries\color{primary}}{\thesection}{0.6em}{}
  [\sectionrule]
\titleformat{\subsection}
  {\large\bfseries\color{secondary}}{\thesubsection}{0.5em}{}
\titleformat{\subsubsection}
  {\normalsize\bfseries}{\thesubsubsection}{0.4em}{}
\titlespacing*{\section}{0pt}{14pt}{6pt}
\titlespacing*{\subsection}{0pt}{10pt}{4pt}

% ---------- Listings (code blocks) ----------
\lstdefinestyle{plainbox}{
  backgroundcolor=\color{codebg},
  basicstyle=\ttfamily\small\color{codefg},
  breaklines=true,
  frame=single,
  framerule=0.4pt,
  rulecolor=\color{rulec},
  xleftmargin=8pt,
  framexleftmargin=8pt,
  framexrightmargin=8pt,
  framextopmargin=4pt,
  framexbottommargin=4pt,
  showstringspaces=false,
  columns=fullflexible,
  keepspaces=true,
}
\lstdefinestyle{python}{style=plainbox,language=Python,
  keywordstyle=\color{primary}\bfseries,
  commentstyle=\color{secondary}\itshape,
  stringstyle=\color{green!50!black},
}
\lstset{style=plainbox}

% ---------- Inline code ----------
\newcommand{\code}[1]{\colorbox{codebg}{\strut\texttt{\small #1}}}

% ---------- Document ----------
\begin{document}

\begin{center}
  {\Huge\bfseries FinTrack Technical Report}\\[0.4em]
  {\large\color{secondary} AI-Powered Finance \& Expense Tracking Platform}\\[0.3em]
  {\small Sagar Makwana \quad\textbar\quad \today}
\end{center}
\vspace{0.5em}
\hrule
\vspace{1em}

% =====================================================================
\section{Project Overview}

FinTrack is an AI-powered finance and expense tracking platform built with a modern full-stack architecture. The application enables users to track income and expenses, set budgets, and receive AI-powered financial insights.

% =====================================================================
\section{Technology Stack}

\subsection{Frontend}
\begin{itemize}[leftmargin=*,itemsep=2pt]
  \item \textbf{Framework}: Next.js 14 (App Router)
  \item \textbf{Language}: TypeScript
  \item \textbf{Styling}: TailwindCSS
  \item \textbf{UI Components}: shadcn/ui (Radix UI primitives)
  \item \textbf{State Management}: TanStack Query (React Query)
  \item \textbf{Charts}: Recharts
  \item \textbf{Notifications}: Sonner (toast)
  \item \textbf{Deployment}: Vercel
\end{itemize}

\subsection{Backend}
\begin{itemize}[leftmargin=*,itemsep=2pt]
  \item \textbf{Framework}: FastAPI
  \item \textbf{Language}: Python 3.11
  \item \textbf{Database}: MongoDB (Atlas)
  \item \textbf{ODM}: Beanie (async MongoDB driver)
  \item \textbf{Authentication}: JWT (access + refresh tokens)
  \item \textbf{AI}: Groq API (llama-3.3-70b-versatile) with OpenAI fallback
  \item \textbf{Logging}: Structlog (JSON logs)
  \item \textbf{Rate Limiting}: slowapi
  \item \textbf{Deployment}: Render/Fly
\end{itemize}

% =====================================================================
\section{Architecture}

\subsection{Design Pattern: Hexagonal (Clean Architecture)}

The backend follows a hexagonal architecture pattern with clear separation of concerns:

\begin{lstlisting}
app/
|-- core/                    # Cross-cutting concerns
|   |-- config.py           # Configuration management
|   |-- logging.py          # Structured logging setup
|   |-- middleware.py       # Request ID middleware
|   |-- rate_limit.py       # Rate limiting
|   |-- security.py         # JWT token operations
|   `-- errors.py           # Domain exceptions
|-- modules/                 # Feature modules
|   |-- auth/
|   |   |-- domain/          # Protocols (repository contracts)
|   |   |-- application/     # Services and DTOs
|   |   |-- infrastructure/  # Beanie models and repositories
|   |   `-- interfaces/http/ # FastAPI routers and deps
|   |-- users/
|   |-- categories/
|   |-- transactions/
|   |-- budgets/
|   |-- analytics/
|   `-- chat/
|-- db.py                    # Database initialization
`-- main.py                  # Composition root
\end{lstlisting}

\subsection{Layer Responsibilities}

\textbf{Core Layer}
\begin{itemize}[leftmargin=*,itemsep=1pt]
  \item Configuration management with Pydantic Settings
  \item Structured JSON logging with request correlation
  \item JWT token generation and validation
  \item Rate limiting per endpoint
  \item Domain exception to HTTP error translation
\end{itemize}

\textbf{Module Layers}
\begin{enumerate}[leftmargin=*,itemsep=1pt]
  \item \textbf{Domain} (\code{domain/}): Defines protocols/interfaces (repository contracts)
  \item \textbf{Application} (\code{application/}): Business logic, DTOs, service orchestration
  \item \textbf{Infrastructure} (\code{infrastructure/}): Beanie ODM models, MongoDB repository implementations
  \item \textbf{Interfaces} (\code{interfaces/http/}): FastAPI routers, dependency injection factories
\end{enumerate}

\subsection{Data Flow}

\begin{lstlisting}
HTTP Request -> Router -> Service -> Repository -> MongoDB
                    |
                    v
                 DTOs (validation)
                    |
                    v
                 Domain Exceptions -> Error Handler -> HTTP Response
\end{lstlisting}

% =====================================================================
\section{Database Schema}

\subsection{Collections}

\textbf{users}
\begin{lstlisting}[style=python]
{
  "email": str (unique),
  "password_hash": str,
  "name": str,
  "role": "user" | "admin",
  "is_active": bool,
  "created_at": datetime
}
Indexes: email (unique)
\end{lstlisting}

\textbf{categories}
\begin{lstlisting}[style=python]
{
  "user_id": ObjectId,
  "name": str,
  "type": "income" | "expense",
  "color": str,
  "created_at": datetime
}
Indexes: user_id+name (unique), user_id+type
\end{lstlisting}

\textbf{transactions}
\begin{lstlisting}[style=python]
{
  "user_id": ObjectId,
  "type": "income" | "expense",
  "amount": float,
  "currency": str,
  "category_id": ObjectId,
  "category_name": str,  # denormalized
  "note": str | null,
  "date": datetime,
  "created_at": datetime
}
Indexes: user_id+date, user_id+category_id+date, user_id+type+date
\end{lstlisting}

\textbf{budgets}
\begin{lstlisting}[style=python]
{
  "user_id": ObjectId,
  "category_id": ObjectId,
  "month": str (YYYY-MM),
  "limit": float,
  "created_at": datetime
}
Indexes: user_id+month+category_id (unique)
\end{lstlisting}

% =====================================================================
\section{API Endpoints}

\subsection{Authentication}
\begin{itemize}[leftmargin=*,itemsep=1pt]
  \item \code{POST /auth/register} --- User registration with default category seeding
  \item \code{POST /auth/login} --- Login with email/password, returns JWT tokens
  \item \code{POST /auth/refresh} --- Refresh access token
  \item \code{POST /auth/logout} --- Logout (client-side token deletion)
\end{itemize}

\subsection{Categories}
\begin{itemize}[leftmargin=*,itemsep=1pt]
  \item \code{GET /categories?type=income|expense} --- List user's categories
  \item \code{POST /categories} --- Create custom category
  \item \code{PATCH /categories/\{id\}} --- Update category
  \item \code{DELETE /categories/\{id\}} --- Delete (rejects if transactions reference it)
\end{itemize}

\subsection{Transactions}
\begin{itemize}[leftmargin=*,itemsep=1pt]
  \item \code{GET /transactions?type=\&category\_id=\&date\_from=\&date\_to=\&limit=\&skip=} --- List with filters and pagination
  \item \code{POST /transactions} --- Create transaction with category validation and denormalization
  \item \code{POST /transactions/suggest-category} --- \textbf{AI-powered category suggestion} from a note + type
  \item \code{GET /transactions/\{id\}} --- Fetch single transaction
  \item \code{PATCH /transactions/\{id\}} --- Update transaction
  \item \code{DELETE /transactions/\{id\}} --- Delete transaction
\end{itemize}

\subsection{Budgets}
\begin{itemize}[leftmargin=*,itemsep=1pt]
  \item \code{GET /budgets?month=YYYY-MM} --- List budgets for month
  \item \code{POST /budgets} --- Upsert by (month, category\_id)
  \item \code{DELETE /budgets/\{id\}} --- Delete budget
\end{itemize}

\subsection{Analytics}
\begin{itemize}[leftmargin=*,itemsep=1pt]
  \item \code{GET /analytics/summary?date\_from=\&date\_to=} --- Total income/expense/net/tx count
  \item \code{GET /analytics/by-category?date\_from=\&date\_to=} --- Breakdown per category
  \item \code{GET /analytics/monthly-trend?months=} --- Last N months income/expense/net
  \item \code{GET /analytics/budget-progress?month=} --- Budget progress per category
\end{itemize}

\subsection{Chat (AI)}
\begin{itemize}[leftmargin=*,itemsep=1pt]
  \item \code{POST /chat} --- Conversational AI with tool-calling for financial data
  \item \code{POST /chat/summary} --- AI-generated financial summary (week/month/quarter/year)
\end{itemize}

\subsection{Meta}
\begin{itemize}[leftmargin=*,itemsep=1pt]
  \item \code{GET /health} --- Application health with version
  \item \code{GET /ready} --- Dependency health (database, AI provider)
\end{itemize}

% =====================================================================
\section{Security}

\subsection{Authentication}
\begin{itemize}[leftmargin=*,itemsep=1pt]
  \item Password hashing with bcrypt (pinned $<$4.1 for compatibility)
  \item JWT access tokens (15min TTL) and refresh tokens (7 days TTL)
  \item Token type validation in JWT payload
  \item Auto-refresh on 401 responses
\end{itemize}

\subsection{Authorization}
\begin{itemize}[leftmargin=*,itemsep=1pt]
  \item All protected routes require valid JWT via \code{get\_current\_user} dependency
  \item User-scoped queries (all operations filtered by user\_id)
  \item Category ownership validation before transaction creation/update
\end{itemize}

\subsection{Rate Limiting}
\begin{itemize}[leftmargin=*,itemsep=1pt]
  \item slowapi with default limits (120/min)
  \item Per-endpoint configuration available
  \item 429 response on rate limit exceeded
\end{itemize}

\subsection{CORS}
\begin{itemize}[leftmargin=*,itemsep=1pt]
  \item Configured whitelist via \code{CORS\_ORIGINS}
  \item Methods: GET, POST, PATCH, DELETE, OPTIONS
  \item Headers: Content-Type, Authorization, X-Request-ID
\end{itemize}

\subsection{Input Validation}
\begin{itemize}[leftmargin=*,itemsep=1pt]
  \item Pydantic models with field validators
  \item Amount limits (transaction: 10M INR, budget: 100M INR)
  \item Currency: 3-char validation
  \item Note: max 500 chars
  \item Month format: YYYY-MM regex validation
\end{itemize}

% =====================================================================
\section{Performance}

\subsection{Database Indexing}
\begin{itemize}[leftmargin=*,itemsep=1pt]
  \item All frequently queried fields have compound indexes
  \item Unique indexes prevent duplicates (email, category name, budget keys)
  \item Date fields indexed descending for time-series queries
\end{itemize}

\subsection{Pagination}
\begin{itemize}[leftmargin=*,itemsep=1pt]
  \item Transaction list: limit (1--200), skip ($\geq 0$)
  \item Prevents large result sets from overloading memory
\end{itemize}

\subsection{Caching}
\begin{itemize}[leftmargin=*,itemsep=1pt]
  \item TanStack Query caches API responses in frontend
  \item Configurable stale time and cache time
\end{itemize}

\subsection{Denormalization}
\begin{itemize}[leftmargin=*,itemsep=1pt]
  \item \code{category\_name} stored in transactions for fast reads without \code{\$lookup}
  \item Trade-off: requires update on category rename (not implemented)
\end{itemize}

% =====================================================================
\section{Observability}

\subsection{Logging}
\begin{itemize}[leftmargin=*,itemsep=1pt]
  \item Structlog with JSON output
  \item Request ID middleware for correlation
  \item Contextual logging in services (user\_id, transaction\_id, etc.)
  \item Log levels: INFO (default), configurable via \code{LOG\_LEVEL}
\end{itemize}

\subsection{Health Checks}
\begin{itemize}[leftmargin=*,itemsep=1pt]
  \item \code{/health} --- Basic health with version
  \item \code{/ready} --- Dependency health (database ping, AI provider check)
\end{itemize}

\subsection{Error Handling}
\begin{itemize}[leftmargin=*,itemsep=1pt]
  \item Domain exceptions (NotFoundError, ValidationError, ConflictError)
  \item Global exception handler translates to HTTP responses
  \item Error codes included in response for client handling
\end{itemize}

% =====================================================================
\section{AI Integration}

\subsection{Tool-Calling}
The chat service uses Groq's (OpenAI-compatible) function calling to access real financial data:

\textbf{Available Tools}
\begin{enumerate}[leftmargin=*,itemsep=1pt]
  \item \code{get\_transactions\_summary} --- Fetch transaction totals for date range
  \item \code{get\_category\_breakdown} --- Spending by category
  \item \code{get\_budget\_progress} --- Budget status for month
\end{enumerate}

\textbf{Flow}
\begin{enumerate}[leftmargin=*,itemsep=1pt]
  \item User sends message via \code{/chat} endpoint
  \item Groq LLM (llama-3.3-70b-versatile) decides which tools to call
  \item Backend executes tools against MongoDB (user-scoped queries)
  \item Tool results returned with INR (\rupee) formatted amounts
  \item AI generates natural language response in INR
\end{enumerate}

\subsection{Provider Selection}
\begin{itemize}[leftmargin=*,itemsep=1pt]
  \item Service prefers \code{GROQ\_API\_KEY} if present (free, fast inference)
  \item Falls back to \code{OPENAI\_API\_KEY} if Groq not configured
  \item Uses AsyncOpenAI client with Groq's OpenAI-compatible endpoint
\end{itemize}

\subsection{Summary Generation}
\begin{itemize}[leftmargin=*,itemsep=1pt]
  \item Period-based summaries (week, month, quarter, year)
  \item \code{date\_from} set to midnight to include first-of-month transactions
  \item Aggregates summary, category breakdown, and budget progress
  \item System prompt enforces INR currency and current date awareness
  \item Returns concise recommendations ($<$200 words)
\end{itemize}

\subsection{AI-Powered Categorization (Bonus)}
\begin{itemize}[leftmargin=*,itemsep=1pt]
  \item Endpoint: \code{POST /transactions/suggest-category}
  \item Fetches the user's categories filtered by transaction type
  \item Sends the note + candidate category names to Groq with temperature 0
  \item Matches LLM response back to a category (exact $\rightarrow$ partial $\rightarrow$ none)
  \item Returns \code{\{category\_id, category\_name, confidence\}} where confidence is \code{high|medium|low}
  \item Frontend exposes this via a ``Suggest'' button beside the Note field; toast confirms the choice
\end{itemize}

\subsection{Real-Time Budget Notifications (Bonus)}
\begin{itemize}[leftmargin=*,itemsep=1pt]
  \item After every expense creation, the frontend re-queries \code{/analytics/budget-progress} for that month
  \item Toast warning (yellow) at $\geq$80\% of budget; toast error (red) at $\geq$100\%
  \item Implemented with Sonner; non-blocking and best-effort (failures swallowed silently)
  \item No websockets required --- leverages existing analytics endpoint for simplicity and zero added infra
\end{itemize}

% =====================================================================
\section{Testing}

\subsection{Backend Tests}
\begin{itemize}[leftmargin=*,itemsep=1pt]
  \item Unit tests for security (JWT hashing, token generation)
  \item Health check tests
  \item Located in \code{tests/} directory
\end{itemize}

\subsection{E2E Tests}
\begin{itemize}[leftmargin=*,itemsep=1pt]
  \item Playwright configuration
  \item Auth flow tests (register, login, redirect)
  \item Transaction flow tests (create, list, delete)
  \item Run with \code{npm run test:e2e} in \code{apps/web}
\end{itemize}

% =====================================================================
\section{Deployment}

\subsection{Backend (Render/Fly)}
\begin{itemize}[leftmargin=*,itemsep=1pt]
  \item Build: \code{pip install -e .}
  \item Start: \code{uvicorn app.main:app --host 0.0.0.0 --port \$PORT}
  \item Environment variables: MongoDB URI, JWT secret, Groq/OpenAI key, CORS origins
\end{itemize}

\subsection{Frontend (Vercel)}
\begin{itemize}[leftmargin=*,itemsep=1pt]
  \item Root directory: \code{apps/web}
  \item Build: \code{npm run build}
  \item Environment variable: \code{NEXT\_PUBLIC\_API\_URL}
\end{itemize}

\subsection{Database}
\begin{itemize}[leftmargin=*,itemsep=1pt]
  \item MongoDB Atlas free tier
  \item IP whitelisting for backend host
  \item Connection string with retry writes
\end{itemize}

% =====================================================================
\section{Product Decisions and Tradeoffs}

\subsection{Why Groq over OpenAI?}
Groq offers an OpenAI-compatible API with the open-source \code{llama-3.3-70b-versatile} model at zero cost and faster inference than gpt-4o-mini. The service is implemented with fallback to OpenAI to keep flexibility.

\subsection{Why MongoDB?}
Financial data is naturally document-oriented (transactions, budgets per user). Beanie ODM provides Pydantic-native models that flow cleanly through validation, business logic, and serialization without translation layers.

\subsection{Denormalization of \code{category\_name}}
Stored on each transaction to avoid \code{\$lookup} joins for list views. Tradeoff: a category rename requires fan-out update (not implemented; acceptable as categories are rarely renamed).

\subsection{Hexagonal Architecture}
Clear separation of \code{domain} / \code{application} / \code{infrastructure} / \code{interfaces} allows testing services with in-memory fakes (see \code{tests/}). Tradeoff: more boilerplate per module, but pays off as the codebase grows.

\subsection{Date Handling in AI Tools}
Tool schemas accept \code{YYYY-MM-DD} (Groq's preferred format). Backend normalizes both \code{YYYY-MM-DD} and ISO 8601. The summary endpoint resets \code{date\_from} to midnight so transactions on the first of the month are not excluded --- a bug discovered and fixed during testing.

% =====================================================================
\section{Challenges Faced}

\begin{enumerate}[leftmargin=*,itemsep=3pt]
  \item \textbf{Groq date format mismatch}: Tool schemas originally specified \code{date-time} format; Groq returned \code{YYYY-MM-DD}. Resolved by accepting both formats in backend parsing.
  \item \textbf{Date filter exclusion bug}: Monthly summary used \code{now.replace(day=1)} which kept the current time, accidentally excluding transactions dated \code{2026-05-01 00:00:00}. Fixed by setting time to midnight.
  \item \textbf{AI currency display}: LLM defaulted to dollars despite system prompt. Fixed by explicitly formatting tool return values with \rupee\ prefix.
  \item \textbf{Double submission on Enter}: Form's native submit + custom \code{onKeyDown} triggered the handler twice. Resolved by relying on form submit only and adding an \code{isSubmitting} guard.
  \item \textbf{Stale closure in summary mutation}: \code{onSuccess} referenced an old \code{conversation} snapshot. Fixed with functional \code{setConversation((prev) => ...)}.
\end{enumerate}

% =====================================================================
\section{Future Scalability Considerations}

\begin{itemize}[leftmargin=*,itemsep=2pt]
  \item \textbf{Read replicas} for MongoDB once query volume grows
  \item \textbf{Redis} for caching analytics aggregations and AI tool results
  \item \textbf{Background jobs} (e.g., Arq/Celery) for periodic summary precomputation
  \item \textbf{CDN} for static assets via Vercel/Cloudflare
  \item \textbf{Sharding} by \code{user\_id} if user count scales horizontally
  \item \textbf{Event sourcing} for transaction audit trail
\end{itemize}

% =====================================================================
\section{Cost Summary (Free Tier)}

\begin{itemize}[leftmargin=*,itemsep=2pt]
  \item \textbf{MongoDB Atlas}: Free (512MB storage)
  \item \textbf{Render}: Free (750 hours/month)
  \item \textbf{Vercel}: Free (100GB bandwidth/month)
  \item \textbf{Groq}: Free tier (generous rate limits)
  \item \textbf{OpenAI} (optional fallback): Pay-as-you-go ($\sim$\$0.15/1M tokens)
\end{itemize}

\textbf{Total: \$0/month with Groq}

% =====================================================================
\section{Conclusion}

FinTrack demonstrates a production-ready full-stack application with:
\begin{itemize}[leftmargin=*,itemsep=2pt]
  \item Clean hexagonal architecture
  \item Comprehensive security (JWT, RBAC, rate limiting, CORS, input validation)
  \item AI-powered insights via Groq with tool calling
  \item \textbf{Bonus}: AI-powered category suggestion + real-time budget threshold notifications
  \item Modern UI with shadcn/ui, dark/light theme, search + filtering
  \item E2E testing with Playwright
  \item Deployment-ready configuration (Vercel + Render + MongoDB Atlas)
\end{itemize}

The application is ready for deployment to production environments.

\end{document}
